% ============================================================
%  BBL Type System  Formal Specification
%  Best Business Language
% ============================================================

\documentclass[11pt, letterpaper]{article}

% --- Core packages ---
\usepackage{amsmath}           % Math environments
\usepackage{amssymb}           % Math symbols (⊥, ⊤, etc.)
\usepackage{stmaryrd}          % Semantic brackets ⟦ ⟧
\usepackage{mathpartir}        % Inference rules (inferrule)
\usepackage{geometry}          % Page layout
\usepackage{hyperref}          % Clickable TOC and cross-refs
\usepackage{xcolor}            % Colors
\usepackage{listings}          % Code listings
\usepackage{titlesec}          % Section formatting
\usepackage{booktabs}          % Nice tables
\usepackage{enumitem}          % List formatting
\usepackage{microtype}         % Better typography

% --- Page geometry ---
\geometry{
  margin=1.25in,
  top=1in,
  bottom=1in
}

% --- Colors ---
\definecolor{bblblue}{RGB}{30, 90, 160}
\definecolor{bblgray}{RGB}{100, 100, 100}
\definecolor{codebg}{RGB}{245, 245, 245}

% --- Hyperref setup ---
\hypersetup{
  colorlinks=true,
  linkcolor=bblblue,
  urlcolor=bblblue,
  citecolor=bblblue,
  pdftitle={BBL Type System},
  pdfauthor={Chris},
}

% --- Section formatting ---
\titleformat{\section}{\large\bfseries\color{bblblue}}{§\thesection}{1em}{}
\titleformat{\subsection}{\normalsize\bfseries}{§\thesection.\thesubsection}{1em}{}

% --- Code listing style (for BBL code examples) ---
\lstdefinelanguage{BBL}{
  morekeywords={
    let, fn, struct, enum, match, if, else, return,
    int, float, decimal, bool, string, void
  },
  sensitive=true,
  morecomment=[l]{//},
  morecomment=[s]{/*}{*/},
  morestring=[b]",
}

\lstset{
  language=BBL,
  backgroundcolor=\color{codebg},
  basicstyle=\ttfamily\small,
  keywordstyle=\bfseries\color{bblblue},
  commentstyle=\itshape\color{bblgray},
  stringstyle=\color{bblblue},
  frame=single,
  framerule=0.4pt,
  rulecolor=\color{bblgray},
  xleftmargin=1em,
  xrightmargin=1em,
  breaklines=true,
  showstringspaces=false,
  tabsize=2,
}

% ============================================================
%  NOTATION MACROS
%  Use these consistently throughout the document.
% ============================================================

% Typing judgment:  Γ ⊢ e : τ
\newcommand{\hastype}[3]{#1 \vdash #2 : #3}

% Type context
\newcommand{\ctx}{\Gamma}

% Empty context
\newcommand{\emptyctx}{\cdot}

% Context extension:  Γ, x : τ
\newcommand{\ctxext}[3]{#1, #2 : #3}

% Metavariables
\newcommand{\tvar}{\alpha}          % Type variable
\newcommand{\tvarf}{\beta}          % Second type variable
\newcommand{\expr}{e}               % Expression metavariable
\newcommand{\typ}{\tau}             % Type metavariable
\newcommand{\typB}{\sigma}          % Second type metavariable

% Built-in types
\newcommand{\tInt}{\mathtt{Int}}
\newcommand{\tFloat}{\mathtt{Float}}
\newcommand{\tDecimal}{\mathtt{Decimal}}
\newcommand{\tBool}{\mathtt{Bool}}
\newcommand{\tString}{\mathtt{String}}
\newcommand{\tVoid}{\mathtt{Void}}
\newcommand{\tBottom}{\bot}

% Type constructors
\newcommand{\tarrow}[2]{#1 \to #2}             % Function type
\newcommand{\tOption}[1]{\mathtt{Option}\langle#1\rangle}
\newcommand{\tResult}[2]{\mathtt{Result}\langle#1,\,#2\rangle}
\newcommand{\tList}[1]{\mathtt{List}\langle#1\rangle}

% Subtyping
\newcommand{\subtype}{\mathrel{<:}}

% Substitution:  [τ / α]
\newcommand{\subst}[2]{[#1 / #2]}

% Semantic brackets
\newcommand{\sem}[1]{\llbracket #1 \rrbracket}

% ============================================================
\begin{document}

% --- Title page ---
\begin{titlepage}
  \centering
  \vspace*{2in}
  {\Huge\bfseries\color{bblblue} BBL Type System}\\[0.5em]
  {\large\color{bblgray} Best Business Language  Formal Specification}\\[2em]
  {\normalsize Version 0.1  Draft}\\[0.5em]
  {\normalsize \today}
  \vfill
  {\small\color{bblgray}
    Mascot: Drizzald \\
    This document is the authoritative formal specification of the BBL type system.
  }
\end{titlepage}

% --- Table of contents ---
\tableofcontents
\newpage

% ============================================================
\section{Introduction}

This document specifies the type system of BBL (Best Business Language).
BBL is a statically typed, business-oriented language with C-style syntax.
Its type system is designed to catch common business-domain errors at compile
time  incorrect currency arithmetic, unchecked absence, and unhandled errors 
while remaining accessible to developers without a formal CS background.

\subsection{Design Goals}

\begin{enumerate}[label=\textbf{G\arabic*.}]
  \item \textbf{Safety.} Ill-typed programs must not compile.
  \item \textbf{Expressiveness.} The type system must model business domain
        concepts (monetary values, nullable fields, fallible operations) without
        requiring boilerplate.
  \item \textbf{Inference.} Type annotations should be optional in most
        local contexts.
  \item \textbf{Readability.} Type errors must be understandable by
        practitioners, not just compiler authors.
\end{enumerate}

\subsection{Notation}

Throughout this document, we use the following conventions.

\begin{center}
\begin{tabular}{ll}
  \toprule
  \textbf{Notation} & \textbf{Meaning} \\
  \midrule
  $\ctx$                          & Type context (maps variables to types) \\
  $\emptyctx$                     & Empty context \\
  $\hastype{\ctx}{e}{\typ}$       & Expression $e$ has type $\typ$ in context $\ctx$ \\
  $\tarrow{\typ}{\typB}$          & Function type from $\typ$ to $\typB$ \\
  $\typ \subtype \typB$           & $\typ$ is a subtype of $\typB$ \\
  $\subst{\typ}{\tvar}$           & Substitute $\typ$ for type variable $\tvar$ \\
  \bottomrule
\end{tabular}
\end{center}

% ============================================================
\section{Primitive Types}

BBL provides the following built-in primitive types.

\begin{center}
\begin{tabular}{lll}
  \toprule
  \textbf{Type} & \textbf{Description} & \textbf{Example literal} \\
  \midrule
  $\tInt$     & 64-bit signed integer            & \lstinline|42| \\
  $\tFloat$   & 64-bit IEEE 754 float            & \lstinline|3.14| \\
  $\tDecimal$ & Arbitrary-precision decimal      & \lstinline|19.99d| \\
  $\tBool$    & Boolean                          & \lstinline|true|, \lstinline|false| \\
  $\tString$  & UTF-8 string                     & \lstinline|"hello"| \\
  $\tVoid$    & Unit  no meaningful value       & (implicit) \\
  $\tBottom$  & Bottom  computation never returns & (implicit) \\
  \bottomrule
\end{tabular}
\end{center}

\subsection{Typing Rules  Literals}

\begin{mathpar}
\inferrule*[Right=T-Int]
  { n \in \mathbb{Z} }
  { \hastype{\ctx}{n}{\tInt} }

\inferrule*[Right=T-Float]
  { f \in \mathbb{R} }
  { \hastype{\ctx}{f}{\tFloat} }

\inferrule*[Right=T-Bool-True]
  { }
  { \hastype{\ctx}{\mathtt{true}}{\tBool} }

\inferrule*[Right=T-Bool-False]
  { }
  { \hastype{\ctx}{\mathtt{false}}{\tBool} }

\inferrule*[Right=T-String]
  { s \in \mathtt{UTF\text{-}8} }
  { \hastype{\ctx}{s}{\tString} }
\end{mathpar}

% ============================================================
\section{Arithmetic}

\subsection{Integer Arithmetic}

\begin{mathpar}
\inferrule*[Right=T-Add-Int]
  { \hastype{\ctx}{e_1}{\tInt} \\ \hastype{\ctx}{e_2}{\tInt} }
  { \hastype{\ctx}{e_1 + e_2}{\tInt} }

\inferrule*[Right=T-Sub-Int]
  { \hastype{\ctx}{e_1}{\tInt} \\ \hastype{\ctx}{e_2}{\tInt} }
  { \hastype{\ctx}{e_1 - e_2}{\tInt} }

\inferrule*[Right=T-Mul-Int]
  { \hastype{\ctx}{e_1}{\tInt} \\ \hastype{\ctx}{e_2}{\tInt} }
  { \hastype{\ctx}{e_1 * e_2}{\tInt} }
\end{mathpar}

\noindent\textbf{Design note:} Integer division is deliberately excluded from
this section pending a decision on whether BBL uses truncating or flooring
semantics. See open question \S\ref{sec:open}.

% ============================================================
\section{Variables and Functions}

\subsection{Variable Lookup}

\begin{mathpar}
\inferrule*[Right=T-Var]
  { (x : \typ) \in \ctx }
  { \hastype{\ctx}{x}{\typ} }
\end{mathpar}

\subsection{Function Types}

A function from $\typ$ to $\typB$ is written $\tarrow{\typ}{\typB}$.

\begin{mathpar}
\inferrule*[Right=T-Abs]
  { \hastype{\ctxext{\ctx}{x}{\typ_1}}{e}{\typ_2} }
  { \hastype{\ctx}{\lambda x{:}\typ_1.\; e}{\tarrow{\typ_1}{\typ_2}} }

\inferrule*[Right=T-App]
  { \hastype{\ctx}{e_1}{\tarrow{\typ_1}{\typ_2}} \\
    \hastype{\ctx}{e_2}{\typ_1} }
  { \hastype{\ctx}{e_1\; e_2}{\typ_2} }
\end{mathpar}

% ============================================================
\section{Composite Types}

\subsection{Structs (Product Types)}

% TODO: define struct typing rules

\subsection{Enums (Sum Types)}

% TODO: define enum / variant typing rules

\subsection{Option}

$\tOption{\typ}$ represents a value that may be absent.
BBL does not have \texttt{null}  absence is always explicit.

% TODO: define Option introduction and elimination rules

\subsection{Result}

$\tResult{\typ}{\typB}$ represents a computation that may fail with error type $\typB$.

% TODO: define Result introduction and elimination rules

% ============================================================
\section{Generics}

% TODO: parametric polymorphism rules

% ============================================================
\section{Subtyping}

% TODO: subtyping relation and rules

% ============================================================
\section{Open Questions}
\label{sec:open}

\begin{enumerate}[label=\textbf{Q\arabic*.}]
  \item \textbf{Integer division semantics.} Truncating (toward zero) or flooring
        (toward negative infinity)?
  \item \textbf{Implicit coercions.} Should $\tInt$ coerce implicitly to
        $\tFloat$? To $\tDecimal$?
  \item \textbf{Variance.} How are generic type parameters treated with respect
        to subtyping  covariant, contravariant, or invariant?
  \item \textbf{Trait / interface bounds.} Syntax and semantics for constraining
        type variables (e.g., \lstinline|T: Numeric|).
\end{enumerate}

% ============================================================
\end{document}